\section{Guidelines}

\section{WUSTL Guidelines}

When you have selected a research topic, you must describe your
research plan as a written dissertation proposal. The proposal, which
typically ranges from 15 to 25 pages in length, should include the
following pieces:

\begin{itemize}

\item motivation for the planned research;

\item  one or more specific aims that you will achieve through your
research program;

\item  background sufficient for someone educated in CSE (but not
necessarily in your area) to understand your specific aims and their
significance;

\item a review of existing work related to your research plan,
indicating how your planned research is novel;

\item high-level details of the planned research, including
challenges, methods, and deliverables, sufficient for an expert in
your area to evaluate its technical feasibility;

\item a timeline indicating when you plan to achieve your specific
aims and when you plan to complete the dissertation.

\item don't forget to describe evaluation of success

\end{itemize}

A good model for the dissertation proposal is an NSF or NIH grant
proposal. Your adviser or another faculty member can provide examples
of such proposals.


\section{Personal Goals}

Goals (what the topic should have):
\begin{itemize}
  \item something theoretical
  \item some application of theory
  \item don't get bogged down in implementation 
  \item sequence of papers that tell a story
  \item a story that is non-big brother
  \item a good job talk... non scary, a good story, a direction of
    research for a few years, shows that I can work with other
    researchers at the institution
\end{itemize}


\section{Meta Proposal Stuff}

\subsection{Evaluating an idea (personal)}
\begin{itemize}
  \item what is the mathy bit?
  \item what is the core theory?
  \item is it a bunch of one-offs?
  \item I want to plant a flag in a real vision problem.
  \item how does it fit into a good job talk
\end{itemize}

What useful things can it do?  Need to tie it back to real
applications.

\subsection{Choosing a Research Problem}

\begin{itemize}
  \item Problem should be important
  \item Problem should hold your personal interest
  \item Problem should have depth, in terms of aspects possibly
    available for investigation
  \item Problem might come from questioning existing literature
  \item Problem may be amenable to some technique you already have devised
  \item Problem may lead you into new technology
  \item Problem may be in ‘hot’ area
\end{itemize}

\subsection{How to proceed?}

\begin{itemize}
  \item Set aside uninterruptible blocks of `research thinking time' in your weekly schedule
  \item Familiarize yourself with previous work from the literature
  \item Critically examine previous approaches, questioning generality, practicality, validation
  \item Frame long-term questions to be answered
  \item Use short-term objectives to subdivide research into manageable pieces
  \item Divide work into investigations that ‘fit’ into a coherent whole
  \item Make progress one paper at a time
  \item Know what it means to ‘solve a problem’ or validate a technique
  \item Write papers and give talks about your work
  \item Develop a personal style -  One at a time vs juggling several projects
\end{itemize}


